\begin{longtable}{|p{0.8cm}|p{11.5cm}|}
	\caption{Kebutuhan Fungsional DecMed \autocite{sadewa2025decmed}}
	\label{tab:kebutuhan-fungsional-decmed}                                                                                                                      \\
	\hline
	\textbf{ID} & \textbf{Deskripsi}                                                                                                                             \\
	\hline
	\endfirsthead

	\multicolumn{2}{c}{\tablename\ \thetable{} -- Lanjutan}                                                                                                      \\
	\hline
	\textbf{ID} & \textbf{Deskripsi}                                                                                                                             \\
	\hline
	\endhead

	\hline
	\endfoot

	F01         & Pasien dapat melakukan registrasi akun pada sistem                                                                                             \\ \hline
	F02         & Pasien dapat masuk ke akun menggunakan credentials yang dimiliki                                                                               \\ \hline
	F03         & Pasien dapat keluar dari akun                                                                                                                  \\ \hline
	F04         & Pasien dapat melakukan scan kode QR yang dimiliki oleh masing-masing tenaga kesehatan dan petugas rekam medis                                  \\ \hline
	F05         & Pasien dapat membaca seluruh RME yang ia miliki                                                                                                \\ \hline
	F06         & Pasien dapat memberikan akses baca bagian administratif dari data RME miliknya kepada petugas rekam medis untuk rentang waktu tertentu         \\ \hline
	F07         & Pasien dapat memberikan akses baca bagian administratif dan klinis dari data RME miliknya kepada tenaga kesehatan untuk rentang waktu tertentu \\ \hline
	F08         & Pasien dapat melihat daftar log dari akses terhadap data RME yang dimiliki                                                                     \\ \hline
	F09         & Pasien dapat mencabut akses yang telah diberikan. Akses tersebut berupa akses baca maupun tulis terhadap data RME.                             \\ \hline
	F10         & Tenaga kesehatan dapat melakukan aktivasi client                                                                                               \\ \hline
	F11         & Tenaga kesehatan dapat melakukan registrasi akun pada sistem                                                                                   \\ \hline
	F12         & Tenaga kesehatan dapat masuk ke akun menggunakan credentials yang dimiliki                                                                     \\ \hline
	F13         & Tenaga kesehatan dapat keluar dari akun                                                                                                        \\ \hline
	F14         & Tenaga kesehatan dapat membaca bagian administratif dan klinis dari data RME milik pasien setelah diberikan akses dalam waktu tertentu         \\ \hline
	F15         & Tenaga kesehatan dapat menambahkan entri RME baru                                                                                              \\ \hline
	F16         & Tenaga kesehatan dapat mengedit bagian klinis dari entri RME baru yang dibuatnya pada rentang waktu yang telah ditentukan                      \\ \hline
	F17         & Petugas rekam medis dapat melakukan aktivasi client                                                                                            \\ \hline
	F18         & Petugas rekam medis dapat melakukan registrasi akun pada sistem                                                                                \\ \hline
	F19         & Petugas rekam medis dapat masuk ke akun menggunakan credentials yang dimiliki                                                                  \\ \hline
	F20         & Petugas rekam medis dapat keluar dari akun                                                                                                     \\ \hline
	F21         & Petugas rekam medis dapat membaca bagian administratif dari data RME milik pasien setelah diberikan akses dalam waktu tertentu                 \\ \hline
	F22         & Admin fasyankes dapat melakukan aktivasi client                                                                                                \\ \hline
	F23         & Admin fasyankes dapat melakukan registrasi akun pada sistem                                                                                    \\ \hline
	F24         & Admin fasyankes dapat masuk ke akun menggunakan credentials yang dimiliki                                                                      \\ \hline
	F25         & Admin fasyankes dapat keluar dari akun                                                                                                         \\ \hline
	F26         & Admin fasyankes dapat menambahkan personel fasyankes baru                                                                                      \\ \hline
	F27         & Admin fasyankes dapat membuat kunci aktivasi baru untuk personel fasyankes                                                                     \\ \hline
	F28         & Kementerian kesehatan dapat mendaftarkan fasyankes dan admin untuk fasyankes tersebut                                                          \\ \hline
	F29         & Kementerian kesehatan dapat membuat kunci aktivasi baru untuk admin fasyankes                                                                  \\ \hline
\end{longtable}